\section*{Week 5}
\section*{Exercise 5.1}
\textbf{Problem:} Use the Fast Powering Algorithm to compute $17^{183} \pmod{256}$.

\textbf{Solution:}

\subsection*{Step 1: Binary Representation of the Exponent}
We express the exponent $183$ in binary form to determine which powers of 2 are needed.
$$ 183_{10} = 128 + 32 + 16 + 4 + 2 + 1 = 10110111_2 $$
\textit{Binary Expansion}

Thus, we can write the base raised to the exponent as:
$$ 17^{183} = 17^{2^7} \cdot 17^{2^5} \cdot 17^{2^4} \cdot 17^{2^2} \cdot 17^{2^1} \cdot 17^{2^0} $$

\subsection*{Step 2: Successive Squaring (Modulo 256)}
We compute the repeated squares of the base $17$ modulo $256$.

\begin{align*}
    17^{2^0} &\equiv 17 \pmod{256} \\
    17^{2^1} &\equiv 17^2 = 289 \equiv 33 \pmod{256} \\
    17^{2^2} &\equiv 33^2 = 1089 = 4 \cdot 256 + 65 \equiv 65 \pmod{256} \\
    17^{2^3} &\equiv 65^2 = 4225 = 16 \cdot 256 + 129 \equiv 129 \pmod{256} \\
    17^{2^4} &\equiv 129^2 = 16641 = 65 \cdot 256 + 1 \equiv 1 \pmod{256} \quad \textit{Cycle Detection}
\end{align*}

Since $17^{2^4} \equiv 1 \pmod{256}$, all subsequent squares will also be $1$:
$$ 17^{2^5} \equiv 1, \quad 17^{2^6} \equiv 1, \quad 17^{2^7} \equiv 1 \pmod{256} $$

\subsection*{Step 3: Product of Coefficients}
Multiply the terms corresponding to the binary bits equal to 1.
\begin{align*}
    17^{183} &\equiv 17^{2^7} \cdot 17^{2^5} \cdot 17^{2^4} \cdot 17^{2^2} \cdot 17^{2^1} \cdot 17^{2^0} \\
    &\equiv 1 \cdot 1 \cdot 1 \cdot 65 \cdot 33 \cdot 17 \pmod{256}
\end{align*}

Calculating the product step-by-step:
\begin{align*}
    17 \cdot 33 &= 561 \equiv 49 \pmod{256} \\
    49 \cdot 65 &= 3185 = 12 \cdot 256 + 113 \equiv 113 \pmod{256}
\end{align*}

\textbf{Final Answer:}
$$ 17^{183} \equiv 113 \pmod{256} $$

\hrule
\vspace{1cm}
\newpage
\section*{Exercise 5.2}
\textbf{Problem:} Find $a \in \{0, \dots, 6\}$ such that $2222^{5555} + 55555^{2222} \equiv a \pmod{7}$.

\textbf{Solution:}

\subsection*{Step 1: Reduce the Bases Modulo 7}
\textit{Linearity of Congruence}
First, we reduce the large bases $2222$ and $55555$ modulo $7$.
\begin{align*}
    2222 &= 317 \cdot 7 + 3 \implies 2222 \equiv 3 \pmod{7} \\
    55555 &= 7936 \cdot 7 + 3 \implies 55555 \equiv 3 \pmod{7} \quad \text{(Note: Student exercise simplified 5555 to 4)}
\end{align*}
\textit{Note regarding source material:} The handwritten notes compute for $5555 \pmod 7 \equiv 4$. Assuming the problem intended $5555$, we proceed with the student's logic: $5555 = 793 \cdot 7 + 4 \equiv 4 \pmod 7$.

Revised Equation based on student notes:
$$ 3^{5555} + 4^{2222} \equiv a \pmod{7} $$

\subsection*{Step 2: Reduce the Exponents using Fermat's Little Theorem}
Since the modulus $p=7$ is prime and $\gcd(3,7)=1$ and $\gcd(4,7)=1$, we apply Fermat's Little Theorem.
\textit{Fermat's Little Theorem: $a^{p-1} \equiv 1 \pmod{p}$}
$$ a^6 \equiv 1 \pmod{7} $$

We reduce the exponents modulo $6$ (since $p-1=6$):
\begin{align*}
    5555 &= 925 \cdot 6 + 5 \implies 5555 \equiv 5 \pmod{6} \\
    2222 &= 370 \cdot 6 + 2 \implies 2222 \equiv 2 \pmod{6}
\end{align*}

Substituting the reduced exponents back into the congruence:
$$ 3^{5555} + 4^{2222} \equiv 3^5 + 4^2 \pmod{7} $$

\subsection*{Step 3: Calculate Final Remainders}
\begin{align*}
    3^5 &= 243 = 34 \cdot 7 + 5 \equiv 5 \pmod{7} \quad (\text{Or } 3^5 = 3^{-1} \equiv 5) \\
    4^2 &= 16 = 2 \cdot 7 + 2 \equiv 2 \pmod{7}
\end{align*}

Adding the results:
$$ 5 + 2 = 7 \equiv 0 \pmod{7} $$

\textbf{Final Answer:}
$$ a = 0 $$

\hrule
\vspace{1cm}
\newpage
\section*{Exercise 5.3}
\textbf{Problem:}
\begin{enumerate}
    \item Compute $2^{11} \pmod{89}$ via the Fast Powering Algorithm.
    \item Deduce the value of $2^{123456789} \pmod{89}$.
\end{enumerate}

\textbf{Solution:}

\subsection*{Part 1: $2^{11} \pmod{89}$ via FPA}
\textit{Binary Representation}
The exponent $11$ in binary is $1011_2 = 8 + 2 + 1$.
$$ 2^{11} = 2^8 \cdot 2^2 \cdot 2^1 $$

Successive squaring modulo 89:
\begin{align*}
    2^{2^0} &= 2 \equiv 2 \pmod{89} \\
    2^{2^1} &= 4 \equiv 4 \pmod{89} \\
    2^{2^2} &= 16 \equiv 16 \pmod{89} \\
    2^{2^3} &= 16^2 = 256 = 2 \cdot 89 + 78 \equiv 78 \pmod{89}
\end{align*}

Multiplying the components:
\begin{align*}
    2^{11} &= 2^8 \cdot 2^2 \cdot 2^1 \\
           &\equiv 78 \cdot 4 \cdot 2 \pmod{89} \\
           &\equiv 78 \cdot 8 \pmod{89} \\
           &= 624 \\
           &= 7 \cdot 89 + 1 \implies 1 \pmod{89}
\end{align*}
\textbf{Result Part 1:} $2^{11} \equiv 1 \pmod{89}$.

\subsection*{Part 2: $2^{123456789} \pmod{89}$ using FLT}
The modulus $89$ is prime. By Fermat's Little Theorem:
\textit{Fermat's Little Theorem}
$$ 2^{88} \equiv 1 \pmod{89} $$

We perform Euclidean division of the exponent $123456789$ by $88$.
$$ 123456789 = 1402918 \cdot 88 + 5 $$

Thus:
\begin{align*}
    2^{123456789} &= 2^{88 \cdot 1402918 + 5} \\
    &= (2^{88})^{1402918} \cdot 2^5 \\
    &\equiv 1^{1402918} \cdot 2^5 \pmod{89} \\
    &\equiv 32 \pmod{89}
\end{align*}

\textbf{Final Answer:}
$$ 32 $$

\hrule
\vspace{1cm}
\newpage
\section*{Exercise 5.4}
\textbf{Problem:} What is the remainder of the division of $221^{654}$ by $109$ (where 109 is prime)?

\textbf{Solution:}

\subsection*{Step 1: Reduce the Base}
\textit{Modular Reduction}
We first reduce $221$ modulo $109$.
$$ 221 = 2 \cdot 109 + 3 \implies 221 \equiv 3 \pmod{109} $$
The problem becomes finding $3^{654} \pmod{109}$.

\subsection*{Step 2: Apply Fermat's Little Theorem}
Since $109$ is prime and $\gcd(3, 109)=1$:
\textit{Fermat's Little Theorem}
$$ 3^{108} \equiv 1 \pmod{109} $$

We divide the exponent $654$ by $108$:
$$ 654 = 6 \cdot 108 + 6 $$

Substituting this into the expression:
\begin{align*}
    3^{654} &= 3^{108 \cdot 6 + 6} \\
    &= (3^{108})^6 \cdot 3^6 \\
    &\equiv 1^6 \cdot 3^6 \pmod{109} \\
    &\equiv 729 \pmod{109}
\end{align*}

\subsection*{Step 3: Final Reduction}
$$ 729 = 6 \cdot 109 + 75 $$

\textbf{Final Answer:}
$$ 75 \pmod{109} $$

\hrule
\vspace{1cm}
\newpage
\section*{Exercise 5.5}
\textbf{Problem:} For $p=47$ and $a=11$, compute $a^{-1} \pmod p$ in two ways:
\begin{enumerate}
    \item[(i)] Extended Euclidean Algorithm (EEA).
    \item[(ii)] Fast Powering Algorithm and Fermat's Little Theorem.
\end{enumerate}

\textbf{Solution:}

\subsection*{(i) Extended Euclidean Algorithm}
We seek $x$ such that $11x \equiv 1 \pmod{47}$, which implies $11x + 47y = 1$.

\textbf{Forward Steps (Euclidean Algorithm):}
\begin{align*}
    47 &= 11 \cdot 4 + 3 \\
    11 &= 3 \cdot 3 + 2 \\
    3 &= 2 \cdot 1 + 1
\end{align*}

\textbf{Backward Substitution (Extended Phase):}
\textit{Bézout's Identity}
Express the remainder 1 as a linear combination of 47 and 11.
\begin{align*}
    1 &= 3 - 1 \cdot 2 \\
      &= 3 - 1 \cdot (11 - 3 \cdot 3) \quad [\text{Substitute } 2] \\
      &= 4 \cdot 3 - 1 \cdot 11 \\
      &= 4 \cdot (47 - 4 \cdot 11) - 1 \cdot 11 \quad [\text{Substitute } 3] \\
      &= 4 \cdot 47 - 16 \cdot 11 - 1 \cdot 11 \\
      &= 4 \cdot 47 - 17 \cdot 11
\end{align*}
Thus, $-17 \cdot 11 \equiv 1 \pmod{47}$.
$$ 11^{-1} \equiv -17 \equiv 30 \pmod{47} $$

\subsection*{(ii) FPA and FLT}
By FLT, since $p=47$ is prime:
$$ 11^{46} \equiv 1 \pmod{47} \implies 11^{-1} \equiv 11^{46-1} \equiv 11^{45} \pmod{47} $$
\textit{Modular Inverse via FLT}

We compute $11^{45}$ using FPA.
Binary of $45$: $45 = 32 + 8 + 4 + 1 = 101101_2$.

\textbf{Squaring:}
\begin{align*}
    11^{2^0} &\equiv 11 \\
    11^{2^1} &= 121 = 2 \cdot 47 + 27 \equiv 27 \equiv -20 \\
    11^{2^2} &= 27^2 = 729 = 15 \cdot 47 + 24 \equiv 24 \\
    11^{2^3} &= 24^2 = 576 = 12 \cdot 47 + 12 \equiv 12 \\
    11^{2^4} &= 12^2 = 144 = 3 \cdot 47 + 3 \equiv 3 \\
    11^{2^5} &= 3^2 = 9
\end{align*}

\textbf{Product:}
$$ 11^{45} = 11^{32} \cdot 11^8 \cdot 11^4 \cdot 11^1 \equiv 9 \cdot 12 \cdot 24 \cdot 11 \pmod{47} $$
Calculation:
\begin{align*}
    12 \cdot 9 &= 108 \equiv 14 \\
    24 \cdot 11 &= 264 \equiv 29 \\
    14 \cdot 29 &= 406 = 8 \cdot 47 + 30 \equiv 30
\end{align*}

\textbf{Final Answer:}
$$ 30 \pmod{47} $$

\hrule
\vspace{1cm}
\newpage
\section*{Exercise 5.6}
\textbf{Problem:} Compute $7^{50} \pmod{47}$.

\textbf{Solution:}

\subsection*{Method: Fast Powering Algorithm}
Exponent $50 = 32 + 16 + 2 = 110010_2$.

\textbf{Successive Squares (Modulo 47):}
\begin{align*}
    7^{2^0} &= 7 \\
    7^{2^1} &= 49 \equiv 2 \\
    7^{2^2} &= 2^2 = 4 \\
    7^{2^3} &= 4^2 = 16 \\
    7^{2^4} &= 16^2 = 256 = 5 \cdot 47 + 21 \equiv 21 \\
    7^{2^5} &= 21^2 = 441 = 9 \cdot 47 + 18 \equiv 18
\end{align*}

\textbf{Calculation:}
$$ 7^{50} = 7^{32} \cdot 7^{16} \cdot 7^2 \equiv 18 \cdot 21 \cdot 4 \pmod{47} $$
\begin{align*}
    18 \cdot 21 &= 378 = 8 \cdot 47 + 2 \equiv 2 \\
    2 \cdot 4 &= 8
\end{align*}

Alternatively, using FLT: $7^{50} = 7^{46} \cdot 7^4 \equiv 1 \cdot 7^4 = 2401 \equiv 4 \pmod{47}$.
\textit{(Note: The student's calculation resulted in 4 in the notes, specifically $2 \cdot 21 \cdot 18 \to 4$. My manual check of $7^4$: $49^2 \equiv 2^2 = 4$. The student likely made a calculation error in the final multiplication step $18 \cdot 21 \cdot 4$. $18 \cdot 21 \equiv 2$. $2 \cdot 4 = 8$. Wait, let me re-verify $7^{2^4}$. $16^2 = 256$. $256 / 47 = 5.44$. $5 \times 47 = 235$. $256 - 235 = 21$. Correct. $7^{2^5} = 21^2 = 441$. $441 / 47 = 9.38$. $9 \times 47 = 423$. $441 - 423 = 18$. Correct. Product: $18 \times 21 \times 4$. $18 \times 21 = 378$. $378 / 47 = 8.04$. $8 \times 47 = 376$. $378 - 376 = 2$. $2 \times 4 = 8$. The FLT method yields 4. Where is the discrepancy? $7^{50} = 7^{46} \cdot 7^4 \equiv 7^4 \pmod{47}$. $7^2=49 \equiv 2$. $7^4 \equiv 2^2 = 4$. The FLT result is definitively 4. In the FPA, $50 = 32 + 16 + 2$. $7^{2} \equiv 2$ (This is $2^1$). Ah, the bit for $2^1$ is 1 ($50 = ...10_2$). $7^{2^1}$ is $2$. $7^{2^4}$ is $21$. $7^{2^5}$ is $18$. $2 \cdot 21 \cdot 18$. $2 \cdot 21 = 42$. $42 \cdot 18 = 756$. $756 / 47 = 16.08$. $16 \times 47 = 752$. $756 - 752 = 4$. Okay, the calculation holds. $18 \cdot 21 \cdot 2 \equiv 4$.)}

\textbf{Final Answer:}
$$ 4 \pmod{47} $$

\hrule
\vspace{1cm}
\newpage
\section*{Exercise 5.7}
\textbf{Problem:} Find all solutions to $x^{1383} \equiv 5 \pmod{13}$.

\textbf{Solution:}

\subsection*{Step 1: Simplify the Exponent}
Since $13$ is prime, we use Fermat's Little Theorem.
\textit{Fermat's Little Theorem Corollary}
For any $x \not\equiv 0$, $x^{12} \equiv 1 \pmod{13}$.
We divide the exponent $1383$ by $12$:
$$ 1383 = 115 \cdot 12 + 3 $$
Thus:
$$ x^{1383} = (x^{12})^{115} \cdot x^3 \equiv 1 \cdot x^3 \equiv x^3 \pmod{13} $$

The equation reduces to finding cubic roots:
$$ x^3 \equiv 5 \pmod{13} $$

\subsection*{Step 2: Trial and Error in $\mathbb{F}_{13}$}
We test values $x \in \{0, \dots, 12\}$ to see which satisfy $x^3 \equiv 5$.

\begin{itemize}
    \item $1^3 = 1$
    \item $2^3 = 8$
    \item $3^3 = 27 \equiv 1$
    \item $4^3 = 64 = 52 + 12 \equiv 12 \equiv -1$
    \item $5^3 = 125 \equiv 8$
    \item $6^3 \equiv (-1) \cdot 8 \equiv 5$ (Wait, $6^3 = 216$. $216 / 13 = 16.6$. $13 \times 16 = 208$. $216-208 = 8$. No.)
    \item Let's check the student's identified solutions: $7, 8, 11$.
    \item Check $x=7$: $7^3 = 343$. $343 / 13 = 26.38$. $26 \times 13 = 338$. $343 - 338 = 5$. \textbf{Solution.}
    \item Check $x=8$: $8^3 = 512$. $512 / 13 = 39.38$. $39 \times 13 = 507$. $512 - 507 = 5$. \textbf{Solution.}
    \item Check $x=11$: $11 \equiv -2$. $(-2)^3 = -8 \equiv 5$. \textbf{Solution.}
\end{itemize}

\textbf{Final Answer:}
The solutions are $x \in \{7, 8, 11\}$.

\hrule
\vspace{1cm}
\newpage
\section*{Exercise 5.10}
\textbf{Problem:} Let $p$ be a prime, $b$ an integer such that $p \nmid b$. Show that $\text{ord}_p(b^{-1}) = \text{ord}_p(b)$.

\textbf{Solution:}

\textit{Definition of Order}
Let $k = \text{ord}_p(b)$. By definition, $k$ is the smallest positive integer such that $b^k \equiv 1 \pmod p$.
Let $h = \text{ord}_p(b^{-1})$. By definition, $h$ is the smallest positive integer such that $(b^{-1})^h \equiv 1 \pmod p$.

\textbf{Step 1: Show $h \mid k$}
Since $b^k \equiv 1 \pmod p$, taking the inverse of both sides:
$$ (b^k)^{-1} \equiv 1^{-1} \implies (b^{-1})^k \equiv 1 \pmod p $$
Since $(b^{-1})^k \equiv 1$, by the property of order, $h$ must divide $k$ ($h \mid k$), so $h \le k$.

\textbf{Step 2: Show $k \mid h$}
Since $(b^{-1})^h \equiv 1 \pmod p$, we can write this as $(b^h)^{-1} \equiv 1$.
Taking the inverse of both sides:
$$ ((b^h)^{-1})^{-1} \equiv 1^{-1} \implies b^h \equiv 1 \pmod p $$
Since $b^h \equiv 1$, by the property of order, $k$ must divide $h$ ($k \mid h$), so $k \le h$.

\textbf{Conclusion:}
Since $h \le k$ and $k \le h$, it must be that $h = k$.
$$ \text{ord}_p(b^{-1}) = \text{ord}_p(b) $$